../cictikz/src/cictikz/data/tex/cictikz_preamble.tex

% The passive mixer: four NMOS switches on the LNA output node, driven by
% the four phase non-overlapping clock.  Each gate is AC coupled to the
% clock and biased through a resistor from V_n.

\begin{circuitikz}[american, thick, transform shape, circuit ee IEC]

  ../cictikz/src/cictikz/data/tex/cictikz_lib.tex

  % one switch: #1 is the y of the switch, #2 the clock name
  \newcommand{\mixsw}[3]{%
    \node[nmos, rotate=-90] (#3) at (0.85,#1) {};
    \draw (#3.S) -- (0,#1);
    \draw (#3.D) -- (2.3,#1);
    \draw (2.42,#1) circle (0.12);
    % the AC coupling capacitor up to the clock port
    \draw (#3.G) -- (0.85,{#1+0.95});
    \draw (0.55,{#1+0.95}) -- (1.15,{#1+0.95});
    \draw (0.55,{#1+1.15}) -- (1.15,{#1+1.15});
    \draw (0.85,{#1+1.15}) -- (0.85,{#1+1.8});
    \draw (0.85,{#1+1.92}) circle (0.12);
    %- beside the port, not above it: stacked above, each label lands
    %- under the device of the switch overhead and reads as its output
    \node[anchor=west] at (1.06,{#1+1.92}) {#2};
    % The bias resistor, tapping the gate side of the capacitor. It has to
    % land below the bottom plate: the whole point of the capacitor is
    % that the gate sits at V_n while the clock swings rail to rail, and
    % a resistor landing between the plates would short that out. Drawn
    % straight rather than diagonally, so it is visible which node it
    % reaches.
    \fill (0.85,{#1+0.72}) circle (0.075);
    \draw (0.85,{#1+0.72}) -- (1.5,{#1+0.72});
    \draw (1.5,{#1+0.72}) to [R] (2.5,{#1+0.72});
    \draw (2.62,{#1+0.72}) circle (0.12);
    \node[anchor=west] at (2.78,{#1+0.72}) {$V_n$};
  }

  \def\yIone{7.2}
  \def\yItwo{4.8}
  \def\yQone{2.4}
  \def\yQtwo{0.0}

  \mixsw{\yIone}{$I_1$}{swione}
  \mixsw{\yItwo}{$I_2$}{switwo}
  \mixsw{\yQone}{$Q_1$}{swqone}
  \mixsw{\yQtwo}{$Q_2$}{swqtwo}

  % the common node, driven by the LNA
  \draw (0,\yQtwo) -- (0,\yIone);
  \draw (0,3.6) -- (-1.3,3.6);
  \draw (-1.42,3.6) circle (0.12);
  \node[anchor=east] at (-1.6,4.0) {LNA};
  \node[anchor=east] at (-1.6,3.45) {Antenna};
  \node[anchor=east] at (-1.6,3.0) {Match};
  \fill (0,3.6) circle (0.075);

  % the I and Q pairs
  \node[anchor=west] at (3.9,6.0) {\Large $I$};
  \node[anchor=west] at (3.9,1.2) {\Large $Q$};

  % the four phase non overlapping clock
  \def\wx{-6.2}      % the time axis origin
  \def\wy{5.2}       % the baseline of the lowest trace
  \draw (\wx,\wy) -- (\wx,{\wy+3.3});
  \draw[->] (\wx,\wy) -- ({\wx+3.9},\wy);

  % Trace #1: baseline y, pulse from x0, one quarter of the period wide.
  % The phases abut rather than leaving a gap, and that is deliberate:
  % four 25% phases fill the period exactly, which is what lets the four
  % switches between them hand off the whole of the LNA current. A dead
  % time here would be signal thrown away, not a safety margin.
  \newcommand{\clktrace}[3]{%
    \draw (\wx,#1) -- ({\wx+#2},#1) -- ({\wx+#2},{#1+0.5})
      -- ({\wx+#2+0.8},{#1+0.5}) -- ({\wx+#2+0.8},#1) -- ({\wx+3.4},#1);
    \node[anchor=east] at ({\wx-0.15},{#1+0.25}) {#3};
  }
  \clktrace{7.9}{0.2}{$I_1$}
  \clktrace{7.0}{1.0}{$Q_1$}
  \clktrace{6.1}{1.8}{$I_2$}
  \clktrace{5.2}{2.6}{$Q_2$}

\end{circuitikz}
\end{document}
